\documentclass[../../../include/open-logic-section]{subfiles}

\begin{document}
\olfileid{sth}{ordinals}{zfm}	
\olsection{$\ZFminus$: એક મહત્ત્વનો પડાવ}

પ્રતિસ્થાપનને કેવી રીતે વાજબી ઠેરવવું, અને ઠેરવી શકાય પણ ખરું કે
નહીં, તે પ્રશ્ન સરળ નથી. તેથી તેની ચર્ચા
\olref[replacement][]{chap} માટે રાખીશું. જોકે પ્રતિસ્થાપન
ઉમેરવાથી આપણે બીજો એક મહત્ત્વનો પડાવ પાર કર્યો છે. હવે સિદ્ધાંત
$\ZFminus$ માટે જરૂરી બધી સ્વયંસિદ્ધિઓ આપણી પાસે છે. વિગતવાર:

\begin{defn}
સિદ્ધાંત $\ZFminus$ની સ્વયંસિદ્ધિઓ આ છે: ઘટકો દ્વારા નિર્ધારિત
સમાનતા, યોગગણ, જોડ, ઘાતગણો, અનંતતા અને પૃથક્કરણ તથા
પ્રતિસ્થાપનના પ્રરૂપોના બધા કિસ્સાઓ. બીજા શબ્દોમાં, $\ZFminus$
એ $\Zminus$માં પ્રતિસ્થાપન ઉમેરે છે.
\end{defn}
\noindent
આ નામ \emph{ઝર્મેલો--ફ્રેન્કેલ} ગણસિદ્ધાંત (કશુંક \emph{બાદ કરેલો};
તેની વાત આગળ આવશે) માટે છે. ફ્રેન્કેલને આ માન મળે છે, કારણ કે
પ્રતિસ્થાપનની રજૂઆત માટે તેમને \citeyear{Fraenkel1922}માં શ્રેય
આપવામાં આવે છે, જોકે તેનું પહેલું ચોક્કસ સ્વરૂપ
\citet{Skolem1922}એ આપ્યું હતું.

\end{document}
